% =========================================================
% EDITE AQUI - RESUMO DO TRABALHO
% =========================================================
% Formato: \siciteresumoitem{Nome do topico}{Texto do topico}
% Regras do template: maximo 2 paginas, sem subsecoes, graficos, figuras
% ou tabelas. Os numeros abaixo sao os obtidos nos experimentos ja
% realizados (CIFAR-10, 5 clientes, 3 rodadas, semente 42).
%
% ATENCAO AO LIMITE DE 2 PAGINAS: este texto tem ~760 palavras, dimensionado
% para caber com folga. Se sobrar espaco apos compilar, os pontos naturais
% para reexpandir sao o Contexto e a lista de trabalhos futuros na Conclusao.

\sicitesection[\sicitespaceAfterResumoTitle]{RESUMO}

\siciteresumoitem{Contexto}{O projeto das redes de comunicação seguiu, por décadas, a formulação de Shannon \cite{shannon1948}, voltada à reprodução fiel de símbolos independentemente de seu significado. Com a saturação do espectro e o crescimento do tráfego gerado por aplicações de inteligência artificial, esse critério tornou-se insuficiente, e as redes 6G apontam para a Comunicação Semântica, em que o objetivo passa a ser o sucesso da tarefa executada no receptor, e não a integridade de cada bit \cite{xie2021deep}. Em paralelo, o Aprendizado Federado permite treinar modelos globais sem centralizar dados \cite{mcmahan2017}, preservando a privacidade dos dispositivos de borda, mas transferindo o gargalo para a comunicação.}

\siciteresumoitem{Lacuna/Problema}{Embora a codificação conjunta fonte-canal profunda esteja consolidada para enlaces ponto a ponto \cite{bourtsoulatze2019deep}, sua integração ao Aprendizado Federado é bem mais recente e ainda pouco consolidada. Falta uma quantificação do compromisso entre o grau de compressão semântica e o desempenho na tarefa quando o treinamento é distribuído entre clientes e a representação transmitida está sujeita a perturbação de canal. Não está estabelecido quanto do tráfego convencional é redundância sem valor para a tarefa, nem qual o custo, em acurácia, de descartá-la.}

\siciteresumoitem{Objetivo}{Quantificar experimentalmente o compromisso entre taxa e desempenho de uma arquitetura federada de comunicação semântica baseada em autoencoders convolucionais leves, medindo a acurácia retida em função da dimensão do espaço latente e do ruído nele injetado, e comparando o custo de comunicação resultante com o da transmissão da imagem íntegra.}

\siciteresumoitem{Método}{A arquitetura é multitarefa. Um codificador convolucional $\mathcal{E}$, embarcável na borda, mapeia a imagem $x$ em um vetor latente $z$ de dimensão $L$; um classificador $\mathcal{C}$ executa a tarefa sobre a representação recebida e um decodificador $\mathcal{D}$ reconstrói a imagem como tarefa auxiliar, regularizando a topologia do espaço latente. A perturbação do enlace é modelada por ruído gaussiano aditivo, $\tilde{z} = z + n$ com $n \sim \mathcal{N}(0,\sigma^{2}I)$, e o treinamento minimiza $\mathcal{L} = \mathcal{L}_{CE}(\mathcal{C}(\tilde{z}),y) + \alpha\,\mathcal{L}_{MSE}(\mathcal{D}(z),x)$, com $\alpha = 0{,}5$. Os clientes treinam localmente e o servidor agrega os parâmetros por \textit{Federated Averaging} ponderado pelo número de amostras \cite{mcmahan2017}. Empregou-se o CIFAR-10 com distribuição independente e idêntica (IID) entre cinco clientes, com três rodadas federadas, uma época local por rodada, otimizador Adam com taxa $10^{-3}$, lotes de 64 amostras e semente fixa, varrendo $L \in \{16, 32, 64, 128\}$ e $\sigma \in \{0;\ 0{,}01;\ 0{,}05;\ 0{,}1\}$. A referência é um classificador convolucional sobre a imagem íntegra, sob o mesmo protocolo federado, com o custo de comunicação contabilizado sobre o número de valores transmitidos por amostra à mesma precisão de 32 bits nos dois casos. A opção por codificadores convolucionais, em lugar de arquiteturas baseadas em atenção usuais na literatura recente de comunicação semântica, é orientada pela restrição de borda: o codificador aqui empregado possui entre $1{,}26 \times 10^{5}$ e $3{,}56 \times 10^{5}$ parâmetros conforme $L$, contra as dezenas de milhões típicas de um \textit{Vision Transformer}.}

\siciteresumoitem{Resultados}{A referência não comprimida atinge acurácia de 0,6771, exigindo $4{,}92 \times 10^{9}$ bits. A configuração com $L=64$ e $\sigma=0{,}05$ atinge 0,6509 com $1{,}02 \times 10^{8}$ bits: redução de 97,9\% no tráfego de inferência em troca de 2,6 pontos percentuais de acurácia. Sob compressão mais severa, $L=16$ preserva 0,6220 com razão de compressão de 192 vezes, equivalente a 99,48\% de economia. A queda de acurácia é fortemente sublinear em relação ao ganho de compressão, sugerindo organização hierárquica da informação no espaço latente: as primeiras dimensões concentram os atributos discriminativos entre classes, enquanto as seguintes acrescentam refinamento pouco relevante para a decisão. Observou-se ainda efeito não monotônico do ruído: para $L=64$, a acurácia sob $\sigma=0{,}05$ superou a do cenário sem perturbação (0,6488), comportamento compatível com a equivalência estabelecida por Bishop \cite{bishop1995noise} entre treinamento com ruído aditivo e regularização de Tikhonov, e com a leitura pela ótica do \textit{Information Bottleneck} \cite{alemi2017vib}. Convém dimensionar esse achado: a diferença é de $2{,}1 \times 10^{-3}$, obtida com uma única semente, não sendo separável da variabilidade de inicialização sem repetição; ademais, $\sigma$ é amplitude absoluta sobre um latente de escala livre, não conversível diretamente em relação sinal-ruído em dB. O resultado é indicativo, não conclusivo.}

\siciteresumoitem{Conclusão}{Os resultados indicam que autoencoders convolucionais leves são candidatos viáveis à codificação semântica na borda, preservando a maior parte do desempenho na tarefa ao transmitir menos de 3\% dos bits originais, o que sustenta a hipótese de que a maior parte do tráfego convencional é redundância sem valor semântico para classificação. Do ponto de vista de projeto de rede, a dimensão latente emerge como parâmetro ajustável de qualidade de serviço: sistemas 6G poderiam negociá-la em função da banda disponível e da criticidade da tarefa, trocando acurácia por taxa de forma controlada e reversível. O ruído do canal, tradicionalmente tratado como degradação, mostra-se aproveitável como regularização obtida gratuitamente do próprio meio. O estudo é exploratório e delimitado pelo escopo avaliado: canal aditivo gaussiano, distribuição homogênea entre clientes, três rodadas federadas, semente única e enlace ideal suposto para a troca de parâmetros. O prosseguimento previsto contempla restrição de potência sobre o latente, que tornaria $\sigma$ interpretável como relação sinal-ruído; desvanecimento Rayleigh e Rician; quantização do latente para ponto fixo; e extensão a cenários não homogêneos com mais rodadas e múltiplas sementes.}
